\chapter{Governing Equations for Frames}\label{ch:frame-equations} 
%Theory of Frames

\section{Introduction}
The \emph{special Cosserat theory of rods} is taken as a starting point, which we regard in the \emph{intrinsic} sense \citep{antman2005nonlinear}\footnote{TODO: Elaborate on intrinsic perspective.}.
This model incorporates the geometrically exact strain measures 
investigated by \cite{cosserat1909theorie}, and can be reduced to
most classical beam models after applying appropriate constraints and assumptions.

In \cref{sec:warping-constraints} warping constraints are introduced in a manner that allows for several different approaches to be unified and implemented in a modular fashion.

A major outcome of this chapter is a foundation for implementing constitutive (section) models that accomodate a wide variety of frame formulations.

\subsection{Literature}

% Differential
%
\begin{itemize}
    \item \cite{euler1744curvis}, \cite{kirchhoff1859ueber}, \cite{love1906treatise}

    \item \cite{cosserat1896theorie}

    \item \cite{ericksen1957exact}, \cite{green1966general},  \cite{naghdi1973theory}.
    From this period, the works of \cite{antman1973theory,antman1974kirchhoffs} were particularly influential. \cite{antman2005nonlinear} is a preeminent source for mathematical treatment of the geometrically exact three-dimensional theory. However, the focus of this work is on analytic results. 

\end{itemize}

% Supplements
\cite{oreilly2017modeling} offers a perspective that is rich in historical context. 
\cite{hjelmstad2005fundamentals} presents a comprehensive treatment of two-dimensional beam theories, with limited coverage of three-dimensional theories and the complications that arise from the non-vectorial manifold.
Useful resources which treat beam theory with  differential geometry are \cite{landau1970theory,connor1976analysis}.

\cite{zienkiewicz2014finite} provides an account of the
necessary theory for its deployment in finite element applications,
while \cite{meier2019geometrically} offers a comprehensive survey of the 
state of the art for geometrically exact finite elements.


\section{Kinematics}


A beam theory is a mathematical model that represents a thin body as a
line embedded in space. This means that every material point in the body
can be identified by a single scalar coordinate \(\xi\).

\subsection{Configuration}

The \emph{configuration} variables for a beam include a vector
  \(\bm{x}\) identifying the centroidal line, a rotation
  \(\FrameRotation\) (small or large) identifying the cross section
  orientation, and a warping amplitude vector
  \(\bm{\alpha}\) which is used to represent deformation modes within a
  cross section.
  \[
  \chi \triangleq \left(\bm{x}, \FrameRotation, \bm{\alpha}\right)
  \]

\subsubsection{Position}

The \emph{position} \(\bm{x}\) is always of size \texttt{ndm} (the dimension of the embedding space).

\subsubsection{Directors}
The \textbf{rotation} variable \(\FrameRotation\) may take several different forms depending on whether small or large deformations are being accomodated, but its \emph{variations} are always a 3-vector for 3D problems and a scalar in 2D problems. 
The rotation \(\FrameRotation\) is conveniently represented as a linear map between
    three director vectors \(\FrameBasis_k\) in the reference
    configuration, and three in the deformed configuration
    \(\frameBasis_k\). The dimension of a rotation variation is denoted
    \texttt{ndr}.

\subsubsection{Warping}
From \cite{epstein2019vlasovs}:
\begin{quote}
In Vlasov’s model, the cross-section can undergo a special kind of warping, whereby small out-of-plane normal displacements can take place. 
This kind of kinematic assumption is shown by \cite{vlasov1961thinwalled} to be suitable for the treatment of the combined bending and torsion of beams of thin-walled open (i.e., simply-connected) cross-section. 
In particular, the planes on which the bending couples act play an important role in the theory in the sense that, when moving a couple to a parallel plane, a moment-couple needs to be taken into consideration. 
This ``moment of a moment'' is what Vlasov calls a \emph{bimoment}.
\end{quote}

\subsection{Deformations}

Objectivity means strains are built from left-invariant differentials of the configuration under the $\mathrm{SE}(3)$ action.
For $\alpha \in W$ the symmetry is translation in $W$ (an Abelian group): the Maurer-Cartan form is just d $\alpha$. Hence the only first-order objective measure is
$$
\alpha^{\prime}, \quad \text { and (optionally) the zero-order field } \alpha \text { itself, }
$$
because an energy or power functional may depend on the "position" in $W$ (hardening/relaxation, penalization, etc.), not just on its gradient.

\subsection{Resultants}


\subsection{Equilibrium}



\section{Director Constraints}

\gls{frame:shear} and \gls{frame:euler}.


\section{Warping Constraints}\label{sec:warping-constraints}

\subsection{Type 1: Independent Warping}

The most general approach to model warping is to take the \(\nwarp\) warping amplitudes \(\bm{\alpha}\) as independent configuration variables. This follows the approach of \citep{reissner1952nonuniform} for warping torsion. Regardless of any kinematic assumptions on the primary kinematic variables \(\FramePosition\) and \(\FrameRotation\), the warping stresses \(\bimoment\) and \(\bishear\) are governed by the same equilibrium equation:
\begin{equation*}
    \bimoment' - \bishear + \bm{f}_{\mathrm{w}} = \bm{0}
\end{equation*}
This can be separated out of the total equilibrium operator \(\operatorname{B}\) of \cref{eq:strong}, so that all assumptions on and simplifications to the primary kinematic variables \(\FramePosition\) and \(\FrameRotation\) are contained within a \emph{primary} equilibrium operator \(\EquilibriumDerivative\). Further investigation of \(\EquilibriumDerivative\) is deffered to \cref{ch:frame-hierarchy}.

A \textbf{Type 1} constitutive formulation is completely general in the sense that it can readily return a response for all other formulation types.

\vspace{1ex}
\begin{ambox}[Type 1: Independent Warping]{box:type1-warping}
A Type 1 constitutive response maps \(6+2\, \nwarp\) strains \((\bm{e}_{\mathrm{p}},\, \bm{\alpha}',\, \bm{\alpha})\) to \(6+ 2\, \nwarp\) generalized stresses \((\bm{s}_{\mathrm{p}},\bm{w},\bm{v})\)
\[
[\bm{s}_{\mathrm{p}},\bm{w},\bm{v}] = \SectionEvaluation{\bm{e}_{\mathrm{p}},\, \bm{\alpha}',\, \bm{\alpha}}
\]
Equilibrium takes the form 
\begin{equation}
\begin{aligned}
    \operatorname{B}_{\mathrm{p}} \bm{\sigma}_{\mathrm{p}} + \bm{f}_{\mathrm{p}} &= \dot{\bm{\mu}} \\ 
    \bimoment' - \bishear + \bm{f}_{\mathrm{w}} &= \bm{0}
\end{aligned}
\end{equation}
\end{ambox}

\subsection{Type 2: Constrained Warping} 

It is common to assume that the amplitudes of certain warping modes are proportional to a component of one of the primary strain variables \(\bm{e}_{\mathrm{p}} = (\ShearStrain,\,\CurveStrain)\). For example, torsional warping is often assumed to be proportional to the rate of twist \(\tau \triangleq \CurveStrain\cdot \FrameBasis\).
Assumptions of this nature can be expressed with \(\mathtt{n_c}\) linear constraints with the form 
\begin{equation}
\mathbf{L}_{\mathrm{w}}\bm{\alpha} - \mathbf{L}_{\mathrm{p}}\bm{e}_{\mathrm{p}} = \bm{0}
\label{eq:linear-warp-constraints}
\end{equation}
where \(\bm{e}_{\mathrm{p}} \triangleq (\ShearStrain, \CurveStrain)\) is the collection of primary strain variables and \(\mathbf{L}_{\mathrm{p}}\) and \(\mathbf{L}_{\mathrm{w}}\) are linear projection operators.

\begin{equation}
    \bm{\alpha} = \bar{\bm{\alpha}} + \mathbf{L}_{\mathrm{p}}\bm{e}_{\mathrm{p}}
\label{eq:decompose-alpha}
\end{equation}
where the reduced warping vector \(\bar{\bm{\alpha}}\)
$$
\bar{\bm{\alpha}} \triangleq \bm{\alpha} - \mathbf{L}_{\mathrm{w}} \bm{\alpha}
$$
has \(\nwarp - \mathtt{n_c}\) non-zero components.


Recall the Galerkin projection \(\mathcal{G}_{\chi}[\bm{\eta}]\) defined by \cref{eq:galerkin}
\begin{equation*}
\mathcal{G}_{\chi}[\bm{\eta}] = \left< \operatorname{B}_{\mathrm{p}} \bm{\sigma}_{\mathrm{p}} ,\, \bm{\eta}_{\mathrm{p}} \right>
+\left< \bm{w}'-\bm{v},\, \delta\bm{\alpha} \right>
\end{equation*}
Using \cref{eq:decompose-alpha} for \(\bm{\alpha}\) and exploiting linearity of the inner product furnishes
\begin{equation*}
\mathcal{G}_{\chi}[\bm{\eta}] = \left< \operatorname{B} \bm{\sigma}_{\mathrm{p}} ,\, \bm{\eta}_{\mathrm{p}} \right>
+\left< \bm{w}'-\bm{v} ,\, \delta \bar{\bm{\alpha}} \right>
+\left< \bm{w}'-\bm{v} ,\, \mathbf{L}_{\mathrm{p}} \delta \bm{e}_{\mathrm{p}} \right>
\end{equation*}
Define the transpose \(\mathbf{L}_{\mathrm{p}}^*\) of \(\mathbf{L}_{\mathrm{p}}\) by
\begin{equation*}
\mathcal{G}_{\chi}[\bm{\eta}] = \left< \operatorname{B} \bm{\sigma}_{\mathrm{p}} ,\, \bm{\eta}_{\mathrm{p}} \right>
+\left< \bm{w}'-\bm{v} ,\, \delta \bar{\bm{\alpha}} \right>
+\left< \mathbf{L}_{\mathrm{p}}^* [\bm{w}'-\bm{v}] ,\, \delta \bm{e}_{\mathrm{p}} \right>
\end{equation*}
\iffalse
By linearity 
$$
\delta \bar{\bm{\alpha}} = \delta \bm{\alpha} - \mathbf{L}_{\mathrm{w}}[\delta \bm{\alpha}]% + \mathbf{L}_{\mathrm{p}}[\delta \ShearStrain, \delta \CurveStrain]
\quad
\text{ and }
\quad
\bar{\bm{\alpha}}' = \bm{\alpha}' - \mathbf{L}_{\mathrm{w}}[\bm{\alpha}']% + \mathbf{L}_{\mathrm{p}}\left[\ShearStrain',\CurveStrain'\right]
$$
\fi 
%
Using \cref{eq:vstrain} for \(\delta \bm{e}_{\mathrm{p}}\triangleq D\bm{e}\cdot \bm{\eta}\) and exploiting linearity of the inner product again furnishes
\begin{equation*}
\mathcal{G}_{\chi}[\bm{\eta}] = 
\left<\bm{s}_{\mathrm{p}} - \mathbf{L}_{\mathrm{p}}^*\bishear + \mathbf{L}_{\mathrm{p}}^* \bimoment' ,\, \delta \bm{e}_{\mathrm{p}} \right>
+\mathcal{G}_{\partial}
+\left< \bm{w}'-\bm{v} ,\, \delta \bar{\bm{\alpha}} \right>
% +\left< \mathbf{L}_{\mathrm{p}}^* [\bm{w}'-\bm{v}] ,\, \delta \bm{e}_{\mathrm{p}} \right>
\end{equation*}
Introduce
\begin{equation}
\begin{aligned}
\bar{\bm{s}}_{\mathrm{p}} &\triangleq \bm{s}_{\mathrm{p}} - \mathbf{L}_{\mathrm{p}}^*\bishear \\
\bar{\bimoment} &\triangleq \bimoment - \mathbf{L}_{\mathrm{w}}^* \bimoment \\
\bar{\bishear} &\triangleq \bishear - \mathbf{L}_{\mathrm{w}}^* \bishear 
\end{aligned}
\label{eq:augment-stress}
\end{equation}
so that
\begin{equation*}
\mathcal{G}_{\chi}[\bm{\eta}] = 
\left< \bar{\bm{s}}_{\mathrm{p}} + \mathbf{L}_{\mathrm{p}}^* \bimoment' ,\, \delta \bm{e}_{\mathrm{p}} \right>
+\mathcal{G}_{\partial}
+\left< \bar{\bm{w}}'-\bar{\bm{v}} ,\, \delta \bar{\bm{\alpha}} \right>
\end{equation*}
Use the relation \(\delta \bm{e}_{\mathrm{p}} = \bm{A}\mathrm{B}^*\bm{\eta}_{\mathrm{p}}\)
and transpose:
\begin{equation}
\mathcal{G}_{\chi}[\bm{\eta}] = 
\left< \operatorname{B}_{\mathrm{p}} \bm{A}_{*} \left[\bar{\bm{s}}_{\mathrm{p}} + \mathbf{L}_{\mathrm{p}}^* \bimoment' \right],\,  \bm{\eta}_{\mathrm{p}} \right>
+\Delta \mathcal{G}_{\partial}
+\left< \bar{\bm{w}}'-\bar{\bm{v}} ,\, \delta \bar{\bm{\alpha}} \right>
\end{equation}
where \(\Delta \mathcal{G}_{\partial} \triangleq\)

\vspace{1ex}
\begin{ambox}[Type 2: Constrained Warping]{box:warp2}
For linear constraints on warping with the form of \cref{eq:linear-warp-constraints}, the material response from Type 1 is augmented with the form
\[
[\bar{\bm{s}}_{\mathrm{p}},\bm{w},\bar{\bm{v}}] = \SectionEvaluation{
\bm{e}_{\mathrm{p}}, \bm{\alpha}', \bar{\bm{\alpha}}
}
\]
with momentum balance given by
\begin{equation}
\begin{aligned}
    \operatorname{B}_{\mathrm{p}} \bm{A}_* \left[\bar{\bm{s}}_{\mathrm{p}} + \mathbf{L}_{\mathrm{p}}^*\bm{w}'\right]+ \bm{f}_{\mathrm{p}} &= \dot{\bm{\mu}} \\ 
    \bar{\bimoment}' - \bar{\bishear} + \bm{f}_{\mathrm{w}} &= \bm{0}
\end{aligned}
\label{eq:type2-equilibrium}
\end{equation}
\end{ambox}

\paragraph{Remarks}
\begin{itemize}
    \item If all warping modes are constrained then the material response has the form
\[
[\bar{\bm{s}}_{\mathrm{p}},\bm{w}] = \SectionEvaluation{\bm{e}_{\mathrm{p}}, \bm{\alpha}'}
\]

\item This condition is well suited for implementation with \gls{frame:euler} models.
\end{itemize}

%

\subsection{Type 3: Uniform Warping}
This case adds to the previous case the condition 
\[
\mathbf{L}_{\mathrm{w}} \bm{\alpha}^{\prime}=\bm{0}
\]
which is consistent with the assumption that warping occurs with a constant amplitude along the length of the body. 

\vspace{1ex}
\begin{ambox}[Type 3: Uniform Warping]{box:warp3}
The augmented constitutive model has the form
\begin{equation*}
[\bar{\bm{s}}_{\mathrm{p}}, \bar{\bimoment},\bar{\bishear}] = \SectionEvaluation{\bm{e}_{\mathrm{p}},\, \bar{\bm{\alpha}}',\, \bar{\bm{\alpha}}
}
\end{equation*}
and momentum balance is prescribed by
\begin{equation}
\begin{aligned}
    \operatorname{B}_{\mathrm{p}} \bm{A}_* \bar{\bm{s}}_{\mathrm{p}} + \bm{f}_{\mathrm{p}} &= \dot{\bm{\mu}}
\end{aligned}
\end{equation}
\end{ambox}

\subsection{Type 4: Neglected Warping}
\[
\mathbf{L}_{\mathrm{w}} \bm{\alpha} =\mathbf{L}_{\mathrm{w}} \bm{\alpha}^{\prime}=\bm{0}
\]
\begin{ambox}[Type 4: Neglected Warping]{box:warp4}
Balance of momentum is stated purely in terms of \(\bm{s}_{\mathrm{p}}\)
\begin{equation}
\begin{aligned}
    \operatorname{B}_{\mathrm{p}} \bm{A}_* \bm{s}_{\mathrm{p}}+ \bm{f}_{\mathrm{p}} &= \dot{\bm{\mu}}
\end{aligned}
\end{equation}
\end{ambox}

\paragraph{Remarks}
\begin{itemize}
    \item This formulation is commonly employed when modeling inelasticity (see, e.g., the \texttt{Fiber} section from OpenSees).
\end{itemize}
% \begin{enumerate}
% \end{enumerate}
